% Осень26!A33; нумерация внутри столбца без пустых строк.
% Источник: лекция 6, §4 (строка 198)
% Источник: лекция 7, §1 (строка 28)
\ticket{30}{Теорема о вычислении интеграла по взвешенному образу меры}
\label{ticket:30}

\materialnote{В источнике доказательство использует теорему Леви и линейность интеграла, доказательства которых отложены.}

\subsection*{Плотность одной меры относительно другой и замена переменных}

Обзорная рамка лектора : прошли сигма-алгебры, полукольца, теорему о продолжении меры, меру Лебега; затем общая теория --- построили интеграл по произвольной мере. Дальше --- ещё общая теория: \textbf{замена переменных в интеграле} (главные приложения будут в \(\mathbb R^m\), где замена переменных --- основной вычислительный приём).

\subsubsection*{Постановка}

\begin{itemize}
\tightlist
\item
  \((X,\mathcal A,\mu)\) --- пространство с мерой (основное).
\item
  \((Y,\mathcal B)\) --- множество с заданной \(\sigma\)-алгеброй \(\mathcal B\) подмножеств \(Y\); \textbf{мера на \(Y\) пока не нужна} (пустое место, «пробельчик»).
\item
  \(\varphi\colon X\to Y\) --- отображение.
\end{itemize}

\subsubsection*{\texorpdfstring{\(\sigma\)-алгебра прообразов; измеримое отображение}{\textbackslash sigma-алгебра прообразов; измеримое отображение}}

\textbf{Факт 1.} Для \textbf{любого} отображения \(\varphi\) (без каких-либо условий на него) система прообразов \[\varphi^{-1}(\mathcal B):=\{\varphi^{-1}(E):\ E\in\mathcal B\}\] --- \(\sigma\)-алгебра в \(X\). \emph{Проверка (упражнение, задано лектором «проверьте, пожалуйста»):} \(\varphi^{-1}(Y)=X\); \(\varphi^{-1}(\complement E)=\complement\,\varphi^{-1}(E)\); \(\varphi^{-1}\bigl(\bigcup_k E_k\bigr)=\bigcup_k\varphi^{-1}(E_k)\) --- и аналогично для пересечений. Всё тривиально проверяется.* Контраст лектора: \textbf{образы} множеств из \(\sigma\)-алгебры при отображении \(\sigma\)-алгебру, вообще говоря, не образуют --- а \textbf{прообразы} «почему-то всегда» образуют.

\begin{definition}

\emph{измеримое отображение.} \(\varphi\colon X\to Y\) \textbf{измеримо}, если \[\varphi^{-1}(\mathcal B)\subset\mathcal A,\] т.е. \(\forall E\in\mathcal B:\ \varphi^{-1}(E)\in\mathcal A\) --- прообраз всякого множества из \(\sigma\)-алгебры \(\mathcal B\) измерим относительно \(\mathcal A\). (Мера \(\mu\) в определении не участвует --- измеримость есть свойство относительно \(\sigma\)-алгебр; «измеримо по мере \(\mu\)» --- синоним в нашей ситуации.)

\end{definition}

\subsubsection*{Образ меры}

\begin{definition}

\emph{образ меры.} Пусть \(\varphi\) измеримо. Определим \(\nu\colon\mathcal B\to[0,+\infty]\) формулой \[\nu(E):=\mu\bigl(\varphi^{-1}(E)\bigr),\qquad E\in\mathcal B.\]

\end{definition}

\textbf{Утверждение.} \(\nu\) --- мера на \(\mathcal B\). Она называется \textbf{образом меры \(\mu\) при отображении \(\varphi\)}.

\emph{Доказательство (счётная аддитивность).} Пусть \(E=\bigsqcup_kE_k\). Тогда \[\varphi^{-1}(E)=\bigsqcup_k\varphi^{-1}(E_k),\] и прообразы дизъюнктны: если точка \(x\) отобразилась в \(E_k\), она уже не может отобразиться в \(E_j\), \(j\ne k\) («\(\varphi\) отправила её в одно место»; склейка точек отображением не мешает --- дизъюнктность прообразов есть общее свойство, инъективность не нужна). Далее по счётной аддитивности \(\mu\): \[\nu(E)=\mu\Bigl(\bigsqcup_k\varphi^{-1}(E_k)\Bigr)=\sum_k\mu\bigl(\varphi^{-1}(E_k)\bigr)=\sum_k\nu(E_k).\]

\textbf{Формула-«кстати» (эквивалентная запись).} Для любого \(E\in\mathcal B\) \[\nu(E)=\int_{\varphi^{-1}(E)} 1\,d\mu\ \Bigl(=\int_X\chi_{\varphi^{-1}(E)}\,d\mu\Bigr)\] --- «доказывать нечего, это и это одно и то же» (свойство 2 из §2).

\subsubsection*{Измеримость композиции}

\textbf{Наблюдение.} Пусть \(f\colon Y\to\overline{\mathbb R}\) измерима относительно \(\sigma\)-алгебры \(\mathcal B\), и \(\varphi\colon X\to Y\) измеримо. Тогда композиция \[f\circ\varphi\colon X\to\overline{\mathbb R},\qquad x\mapsto f(\varphi(x)),\] измерима относительно \(\mathcal A\).

\begin{proof}

Для произвольного \(a\in\mathbb R\) рассмотрим типичное множество Лебега: \[\{x\in X:\ f(\varphi(x))<a\}=\varphi^{-1}\bigl(\{y\in Y:\ f(y)<a\}\bigr)=\varphi^{-1}\bigl(Y(f<a)\bigr).\] Множество \(Y(f<a)\) лежит в \(\mathcal B\) (измеримость \(f\)), поэтому его прообраз лежит в \(\mathcal A\) (измеримость \(\varphi\)).

\end{proof}

Это «замечательный задел» для замены переменных: подставляя вместо переменной \(y\) функцию \(\varphi(x)\), мы из измеримой функции получаем измеримую: \(f(y)\ \leadsto\ f(\varphi(x))\).

\subsubsection*{Взвешенный образ меры}

\begin{definition}

\emph{вес; взвешенный образ меры.} Пусть \(\omega\colon X\to[0,+\infty]\) --- неотрицательная измеримая функция на \(X\) (\textbf{вес}, весовая функция), \(\varphi\colon X\to Y\) измеримо. Определим \(\nu\colon\mathcal B\to[0,+\infty]\): \[\nu(E):=\int_{\varphi^{-1}(E)}\omega\,d\mu\ =\ \int_X\omega\,\chi_{\varphi^{-1}(E)}\,d\mu,\qquad E\in\mathcal B.\]

\end{definition}

\textbf{Утверждение.} \(\nu\) --- мера на \(\mathcal B\); она называется \textbf{взвешенным (с весом \(\omega\)) образом меры \(\mu\)}. \emph{Обоснование --- «по аналогии с предыдущим, ничего не изменится»:} счётная аддитивность \(\nu\) --- это в точности счётная аддитивность интеграла по множеству (§3), применённая к неотрицательной измеримой \(\omega\): \(\nu(\sqcup E_k)=\int_{\sqcup_k\varphi^{-1}(E_k)}\omega\,d\mu=\sum_k\int_{\varphi^{-1}(E_k)}\omega\,d\mu\). (Лектор свернул проверку; при \(\omega\equiv1\) получаем невзвешенный образ из п. 4.3.)

\subsubsection*{Теорема о замене переменных («зоологическая теорема»)}

\begin{theorem}

Пусть \((X,\mathcal A,\mu)\) --- пространство с мерой, \((Y,\mathcal B)\) --- множество с \(\sigma\)-алгеброй, \(\varphi\colon X\to Y\) измеримо, \(\nu\) --- взвешенный с весом \(\omega\ge0\) (измеримая) образ меры \(\mu\) при \(\varphi\). Пусть \(f\colon Y\to[0,+\infty]\) --- неотрицательная измеримая функция. Тогда \(f\circ\varphi\) измерима на \(X\) (наблюдение 4.4), произведение \(f(\varphi(x))\,\omega(x)\) неотрицательно измеримо, и \[\boxed{\ \int_Y f(y)\,d\nu(y)\ =\ \int_X f\bigl(\varphi(x)\bigr)\,\omega(x)\,d\mu(x)\ }\qquad(\star)\] \emph{(«Зоологическая» --- по лектору: ничего не придумывается, «садимся и наблюдаем», как определение интеграла работает при пересадке мер.)}

\end{theorem}

\textbf{Обозначение} \(d\nu(y)\), \(d\mu(x)\): буква переменной под \(d\) лишь акцентирует, в каком пространстве живёт интегрирование (\(y\) бегает по \(Y\), \(x\) --- по \(X\)); с самого курса использовался короткий знак \(\int f\,d\mu\).

\subsubsection*{Доказательство}

\textbf{Шаг 0 (характеристическая функция).} \(f=\chi_B\), \(B\in\mathcal B\). Левая часть: \(\int_Y\chi_B\,d\nu=\nu(B)\) (свойство 2). Правая часть: \(\chi_B(\varphi(x))=\chi_{\varphi^{-1}(B)}(x)\) --- единица стоит ровно на прообразе: \[\int_X\chi_B\bigl(\varphi(x)\bigr)\,\omega(x)\,d\mu=\int_X\chi_{\varphi^{-1}(B)}\,\omega\,d\mu=\int_{\varphi^{-1}(B)}\omega\,d\mu=\nu(B)\] --- по определению интеграла по множеству (характеристическая функция «вырезает» множество интегрирования) и определению взвешенного образа. Обе части совпали. \ensuremath{\checkmark}

\textbf{Шаг 1 (ступенчатые).} \(f=\sum_{\text{кон}}\lambda_i\chi_{B_i}\). Формула \((\star)\) следует из шага 0 и \textbf{линейности интеграла по функции}. \textbf{Требует сверки.} \textbf{Лектор признал: аддитивная часть линейности --- \(\int(f+g)=\int f+\int g\) --- к этому моменту ещё НЕ доказана} (однородность \(\int(\alpha f)=\alpha\int f\) доказана сегодня, свойство 4; порядок изложения был «потороплен» --- «я вообще зря вам эту теорему рассказал» ). Будет доказана на следующей лекции.

\textbf{Шаг 2 (неотрицательная измеримая).} Интеграл \(\int f\,d\nu\) --- супремум интегралов ступенчатых, но супремум можно \textbf{реализовать вдоль последовательности}: по теореме об аппроксимации (характеризация неотрицательных измеримых функций ступенчатыми) существует последовательность ступенчатых \(g_n\): \[0\le g_n\le g_{n+1},\qquad g_n(x)\to f(x)\ \ \forall x.\] Для каждой \(g_n\) формула \((\star)\) верна (шаг 1): \[\int_Y g_n\,d\nu=\int_X g_n\bigl(\varphi(x)\bigr)\,\omega(x)\,d\mu.\] Хотим в обеих частях перейти к пределу при \(n\to\infty\) (\textbf{предельный переход под знаком интеграла}). Он закончен, поскольку последовательности монотонны:

\begin{itemize}
\tightlist
\item
  \(g_n\uparrow f\) \ensuremath{\Longrightarrow} \(\int g_n\,d\nu\to\int f\,d\nu\);
\item
  \(g_n(\varphi)\,\omega\uparrow f(\varphi)\,\omega\) (множитель \(\omega\ge0\)) \ensuremath{\Longrightarrow} \(\int g_n(\varphi)\omega\,d\mu\to\int f(\varphi)\omega\,d\mu\). Это --- \textbf{в чистом виде условия теоремы Беппо Леви} (монотонной сходимости): ступенчатость не нужна, достаточно неотрицательности и монотонности. \textbf{Требует сверки.} \textbf{Теорема Леви на момент лекции не доказана} --- объявлена первоочередной теоремой следующей лекции («смотри доказательство на следующей странице, условно говоря»). Формально шаг 2 завершится после неё.
\end{itemize}

\textbf{Дополнение (суммируемый случай).} Та же формула \((\star)\) верна, если вместо неотрицательной измеримой \(f\) взять \(f\), \textbf{суммируемую по мере \(\nu\)}. \emph{Обоснование.} \(|f|\ge0\) измерима \ensuremath{\Longrightarrow} по уже полученному (шаги 0--2): \[\int_Y|f|\,d\nu=\int_X\bigl|f\bigl(\varphi(x)\bigr)\bigr|\,\omega(x)\,d\mu<\infty,\] и правая часть конечна. Так как \(\omega\ge0\), имеем \(|f(\varphi)\,\omega|=|f(\varphi)|\,\omega\), поэтому функция \(f(\varphi)\,\omega\) суммируема по мере \(\mu\) (критерий суммируемости, §1). Раскладываем \(f=f^+-f^-\): к \(f^\pm\) (неотрицательным измеримым) формула \((\star)\) применима, обе части --- конечные величины, и их можно вычитать: \[\int_Y f\,d\nu=\int_Y f^+\,d\nu-\int_Y f^-\,d\nu=\int_X f^+(\varphi)\omega\,d\mu-\int_X f^-(\varphi)\omega\,d\mu=\int_X f(\varphi)\,\omega\,d\mu. \qquad\blacksquare\]

\textbf{Итог лектора по теореме .} В доказательстве «повисли два должка»: \textbf{теорема Леви} и \textbf{линейность (аддитивность) интеграла} --- после них текст становится осмысленным. План следующей лекции: сначала Леви (предельный переход под знаком интеграла --- ещё до линейности!), затем линейность. Общая мораль: утверждения об интегралах достаточно доказывать для ступенчатых функций, а затем «поднимать» их на измеримые функции монотонным предельным переходом по теореме Леви --- «предел более понятная операция, с которой удобнее пользоваться», чем переход к супремуму, которым лектор обходился до сих пор.

\subsection*{Следствие теоремы о замене переменных: интегрирование по множеству}

\textbf{Напоминание (теорема об интегрировании по взвешенному образу меры --- лекция 6, §4.6).} Пусть \((X,\mathcal A,\mu)\) --- пространство с мерой, \((Y,\mathcal B)\) --- множество с \(\sigma\)-алгеброй, \(\varphi\colon X\to Y\) измеримо, \(\omega\colon X\to[0,+\infty]\) --- измеримый вес, и \(\nu\) --- взвешенный образ меры \(\mu\) при \(\varphi\): \[\nu(B)=\int_{\varphi^{-1}(B)}\omega\,d\mu\qquad(B\in\mathcal B).\] Тогда для всякой суммируемой (или неотрицательной измеримой) \(f\colon Y\to\overline{\mathbb R}\): \[\int_Y f\,d\nu=\int_X f\bigl(\varphi(x)\bigr)\,\omega(x)\,d\mu(x).\]

\textbf{Следствие (интегрирование по части пространства).} В тех же условиях: если \(f\) суммируема на \(Y\) (по мере \(\nu\)) и \(B\in\mathcal B\), то \[\boxed{\ \int_B f(y)\,d\nu(y)\ =\ \int_{\varphi^{-1}(B)} f\bigl(\varphi(x)\bigr)\,\omega(x)\,d\mu(x)\ }\] --- «вместо интегрирования по \(B\) можно сделать замену переменных \(y=\varphi(x)\) и интегрировать по прообразу \(B\) композицию с \(\varphi\), умноженную на вес \(\omega\), по мере \(\mu\)».

\begin{proof}

Применить теорему к функции \(f\cdot\chi_B\). Слева \(\int_Y f\chi_B\,d\nu=\int_B f\,d\nu\) (множитель \(\chi_B\) «вырезает» \(B\)). Справа \(\chi_B(\varphi(x))=\chi_{\varphi^{-1}(B)}(x)\), поэтому \[\int_X f(\varphi(x))\,\chi_{\varphi^{-1}(B)}(x)\,\omega(x)\,d\mu(x)=\int_{\varphi^{-1}(B)}f(\varphi(x))\,\omega(x)\,d\mu(x).\]

\end{proof}

\textbf{Замечание лектора.} Теорема доказывалась по схеме последовательного усложнения подынтегральной функции: «нулевой» шаг --- заготовка проще ступенчатой (одна ступенька = одно множество), затем ступенчатые, неотрицательные измеримые, суммируемые. Именно в «нулевом» шаге видно, как при подстановке \(f\chi_B\) «сама собой» возникает характеристическая функция прообраза \(\chi_{\varphi^{-1}(B)}\): слева это очевидно, справа «это зарыто в доказательстве». \textbf{Требует сверки.} (формальный статус «нулевого» шага восстановлен по смыслу.)
